\summarychapter{45}{The Medicine of Mercy}{Roman Catholicism: Vatican II}

\begin{inbrief}
The Second Vatican Council (1962--1965) replaced the picture of the Roman Church as a besieged
fortress with that of a ``Pilgrim People'' moving through the world. The chapter tells how Pope John
XXIII called the council to bring the church up to date, how conservatives and progressives
contended in its four sessions, and how its work released a wave of change and a crisis of authority
and identity in the decade that followed.
\question{How can we explain this revolution in Rome? What challenged the besieged mentality of
Roman Catholicism?}
\end{inbrief}

Vatican II was a spectacle. At the four-hour opening the white miters of the bishops rippled through
St.~Peter's Square, the pope's portable throne bobbed above them like a royal barge between two
ostrich-feather fans, and the choir sang \emph{Ubi caritas et amor, Deus ibi est}, ``Where charity
and love are found, God is present.'' But it was far more than a show. It revealed a new spirit in
Catholicism crying out for change, shattered the Protestant picture of a monolithic and absolutist
system, and marked the church's first tacit admission that those who left it in the past may have
had good cause. The traditional image, shaped largely by Trent, was of an impregnable fortress
attacked by secularism, modernism and individualism, within whose walls were the changeless truths of
God, the true sacrifice of the Mass and papal infallibility. Pius XII (1939--1958) reflected it in
1950 when the encyclical \emph{Humani Generis} condemned theologians who tried to free the dogmas
from long-held formulations and return to the language of the Bible and the Fathers. Vatican II spoke
instead of pilgrims caring for the weak and weary. It claimed few successes or certainties, for
pilgrims are still on the way.

\ssection{The Good Shepherd, Pope John}
The man who threw open the windows was Angelo Roncalli, John XXIII. He often said that the idea of a
council came to him as a sudden inspiration of the Holy Spirit. Elected on October 28, 1958, he
announced his plan on January 25, 1959, after barely ninety days. Yet he was singularly prepared, an
almost intuitive judge of human hopes and needs. As a young priest he taught the church fathers for
a year at the Pontifical Lateran Seminary until his superiors decided he was not quite safe, since he
suggested, for instance, that mixed marriages might sometimes be allowed. He copied letters in the
Vatican's Oriental Congregation until an apostolic visitor was needed in remote Bulgaria (1925--34),
then spent ten years in Muslim Turkey, and was sent to troubled France near the end of the war only
because Rome would not spare a top man for so messy a post. His humility and ability charmed the
French. In 1953 Pius XII made him a cardinal and patriarch of Venice.

Everywhere he made friends with non-Catholics. In Turkey he helped to rescue and provide for Jews
fleeing Nazi Germany, and in France he recoiled from the films of the bodies at Buchenwald and
Auschwitz: ``How could this be? The mystical body of Christ!'' To a group of Jewish visitors he
simply repeated the biblical greeting, ``I am Joseph, your brother.'' As pope he asked to be known
not as a political or learned pope but as ``the good shepherd defending truth and goodness.'' He
went out to orphanages, jails, schools and churches, once received a traveling circus and patted a
lion cub named Dolly, and ended the custom of clearing visitors from St.~Peter's dome while the pope
walked in the garden below: ``Why shouldn't they look? I'm not doing anything scandalous.''

The council's purpose, he said, was \emph{aggiornamento}, ``bringing up to date,'' a word that
implies not only adaptation to the outward life of the time but an inward change of thought. He
apparently meant the council to turn from the rigid legal patterns of the past to the pastoral
concerns of the present. Many expected a transformation as great as Constantine's conversion or the
Reformation: the overthrow of the tie between religion and politics that Constantine had introduced
and the renunciation of the narrow spirit bred by the Counter-Reformation. Certainly he wanted a
church in tune with the contemporary world.

In October 1962, 2,540 of the 2,908 eligible cardinals, patriarchs, bishops and abbots came to Rome.
Vatican I had had six or seven hundred fathers, and only about two hundred voted on the decrees of
Trent. Now there were 430 Italians, 230 from the United States, which had been mission territory
until 1908, 230 Africans and more than 300 Asians. It was also the first council not called to combat
heresy, define new dogma or marshal the church against enemies. Trent had upheld indulgences and
built a doctrinal fortress against Protestantism. John showed that this council was for something.
His opening speech was a mandate for a pastoral council. The days of the state church are over, and
in the age being born the church must not keep its authority by weapons of repression but ``rule
with the medicine of mercy rather than with severity.'' Its renewal would bring nearer the reunion of
Christendom, when Christ's prayer ``that they may be one'' would be fulfilled.

These brave words were tested two days later at the first working session, when the fathers were to
elect ten permanent commissions. On their desks, for ``guidance,'' lay lists of the members of the
preparatory commissions, which had been directed by cardinals of the Curia, the Italian-controlled
civil service of the church in the Vatican. The suggestion was that all would go smoothly if the
Curia's experts were elected. The issue was who would run the council, the Curia or the fathers, and
behind it how open the Vatican would be to the outside world. Before any vote Cardinal Achille
Li\'enart, the seventy-eight-year-old archbishop of Lille, rose from the presidency table: ``We will
not accept the lists of candidates prepared for us before the Council opened,'' and he asked for a
delay so that the fathers could choose their own. The motion was seconded and approved, and the
session ended in less than half an hour. Li\'enart, who had prepared for six months, saw a maneuver
to give the council a particular direction where others saw a harmless way of saving time. When the
commissions were elected two days later the list was balanced and international.

\ssection{Conservatives and Progressives}
This drew the lines between two groups of bishops. The conservatives, said the Jesuit Bible scholar
Francis J.~McCool, were those ``who see in the future a threat to the past,'' and the progressives
those ``who see in the future the promise of the future.'' They collided on almost every issue. The
conservative position was chiefly that of the Curia. Its members, mostly aging Italians insulated
from the modern world, exerted vast influence over world Catholicism and over the pope himself, were
generally satisfied with the church as it was, and met proposals for change with deep hostility.
Their representative was the erudite Cardinal Alfredo Ottaviani, seventy-two, born in his father's
bakery in the Trastevere quarter of Rome, who had spent most of his life within the one square mile
of the Vatican. He was in a sense the prisoner of his function, which was to watch over the purity
of doctrine and condemn heresy. His name appeared twenty-three times in the pontifical yearbook: he
sat on seven congregations, two commissions and a tribunal and was protector of twenty-two religious
orders. Critics said he thought everything new wrong; what he actually said was that everything new
is not necessarily right. Many such officials had not wanted the council at all. They called it in
private ``the Pope's folly,'' tried to sabotage and delay the preparations for fear it would expose
the church's internal differences, and were determined to keep the Curia's powers of decision in
Italian hands. The progressives had no single leader but gathered around able cardinals from Germany,
France, Austria and the Low Countries, such as the tall, athletic Jan Alfrink of Holland, whose many
contacts with the Protestant ``separated brethren'' led an Italian newspaper to brand him
``anti-Roman.''

The council worked in four sessions, the first from October 11 to December 8, 1962, and three more
in the autumns of the next three years. In the first the progressives wanted the liturgy changed to
allow modern languages in place of Latin and to encourage lay participation in the Mass, and the
conservatives objected. A deeper issue arose over the draft (\emph{schema}) on divine revelation
prepared by Ottaviani's theological commission, which insisted on the two sources of revelation,
Scripture and tradition, recognized since Trent. The progressives saw no point in stressing
differences with Protestants and wished to present the two as channels of a single stream. Are some
truths of the faith found only in tradition, or all of them in Scripture? After nearly two weeks
1,368 fathers voted to shelve the draft, short of the two-thirds required. Pope John, watching on
closed-circuit television, intervened and had it rewritten by a new committee chaired jointly by
Ottaviani and Cardinal Augustin Bea, the Jesuit head of the new Secretariat for Promoting Christian
Unity and a leader of the progressives. ``This day,'' said the Canadian theologian Gregory Baum,
``will go down in history as the end of the Counter Reformation.''

On June 3, 1963, amid preparations for the next session, John XXIII died, and the whole religious
world mourned. On June 21 Cardinal Montini, archbishop of Milan, became Paul VI and at once announced
that the council would continue. In the second session (autumn 1963) the two understandings of the
church clashed. Vatican I had proclaimed the pope's primacy and infallibility; Vatican II tried to
explain how the whole body, or college, of bishops is related to the pope in governing the church.
Conservatives defended papal sovereignty, and progressives pressed for greater power for the bishops
(collegiality). The third session (September 14 to November 21, 1964) debated religious liberty:
would a declaration make God's truth relative and feed indifferentism, the notion that it does not
matter what one believes so long as one is sincere? It also dealt with the life and ministry of
priests, the mission of laymen in the world and missionary work in non-Christian areas.

The last session (September 14 to December 8, 1965) returned to freedom of religion. The declaration
on freedom of conscience proclaimed that no state has the right to prevent by outward pressure the
preaching and acceptance of the gospel. At the same time the church gave up the assumption, held
since Constantine, that where it has the means, as in Spain and Italy, it may use public power to
enforce its religious demands and further its saving work. In accepting the decree, proclaimed on
December 7, Roman Catholicism solemnly renounced in principle the use of external force against
conscience, a radical break with fifteen hundred years of practice. A celebration in St.~Peter's
Square closed the council on December 8, 1965. ``Revolutionary'' is probably too strong a word for
it, Shelley judges. The traditional theology and the papal government remained intact, and the
sixteen decrees were, with few exceptions, compromises between the two parties.

\ssection{A Tidal Wave of Change}
Still, the council broke with the angry spirit of Trent and the defensive mood of Vatican I and
turned the face of Catholicism toward the world in concern, not anger. Its work set off a tidal wave.
The following decade was the most tumultuous in the church's modern history, and so many landmarks
were swept away that the ordinary Catholic was left bewildered. The first wave was the new liturgy.
Taught to regard the Mass as a mysterious and unchangeable ceremony going back to Christ himself, he
now saw the altar brought forward and the priest facing the people and reading the prayers aloud in
their own language where he had whispered them in Latin. Many old ceremonies were dropped, and
worshipers who had kept to their private prayers were asked to greet one another with a ``sign of
peace.''

The most fundamental issue was authority. Until Pope John and the council the typical Catholic took
the authoritarian structure of the church as part of divine revelation, thought of the pope as a
superhuman ruler whose every word had supernatural authority, and regarded even the bishop with awe.
Few questioned procedures that looked medieval to outsiders, and no one challenged the bishop who
ruled his diocese as a personal fief or the pastor who ran his parish. The council raised the
progressives' hopes, but they soon found that little changed in practice. The structure remained a
pyramid, with power flowing down from an infallible head. The council had taken care to safeguard
his absolute authority and gave the bishops only a consulting role, which the pope was free to use
or not. In the light of the council's own new understanding of the church many found this
intolerable and agitated for democratic reforms with dissent, demonstrations, sit-ins and the press.

Tension reached its height in 1968, when Paul VI issued the encyclical \emph{Humanae Vitae}
condemning artificial methods of contraception. He staked his authority on a decision taken against
the overwhelming majority of his own birth control commission, and the affair brought on one of the
gravest crises of papal authority since Luther. A strong body of opinion blamed him for acting
without the bishops. Its chief spokesman was Cardinal Suenens, archbishop of Malines in Belgium and
one of the architects of the council, who called in speeches, press conferences and writings for an
end to the medieval papacy: the pope should no longer act as though he were outside or above the
church. Almost as serious was the prohibition of divorce. By church law no truly sacramental marriage
between baptized Catholics can be dissolved, even by the pope; partners who can no longer live
together may be permitted to separate but not to remarry while either lives. Despite the hardship
few had dared to challenge the law. Once the facade of immutability began to crack, priests and
theologians questioned its wisdom and its scriptural basis. They wanted the church to go on
preaching marriage as a sacrament and a lifelong commitment, and asked for a more flexible pastoral
approach to couples whose marriages had failed.

\ssection{The Crisis of Identity}
Amid the upheaval priests, brothers and nuns left in large numbers. Between 1962 and 1974 the number
of seminarians in the United States fell by 31 percent, and between 1966 and 1972 nearly eight
thousand American priests left the public ministry. The root seemed to be the priest's identity. The
sacredness that once surrounded the priesthood looked odd beside the New Testament's descriptions of
a minister, the democratic trend made the old priestly caste seem medieval, and as church membership
in a secular society became a matter of personal commitment and not of being born Catholic, the
priest no longer needed to carry the whole burden of the church's mission. Suenens traced the
conflict over authority to two theologies. One sees the church above all as a fellowship of
spiritual communities held together in essentials by their recognition of papal primacy. The other,
the traditional one, sees it as a medieval superstate under an absolute monarch who imposes his will
on its members and on society.

Suenens found hope in a new movement of personal Christianity, the Catholic Charismatic Renewal. Its
leaders trace it to the spring of 1966, when two laymen on the faculty of Duquesne University in
Pittsburgh realized that they lacked the early Christians' power to proclaim the gospel, gave
themselves to prayer and shared their concern with colleagues. In August two young men at the
National Cursillo Convention (the Cursillo being a Catholic renewal movement born in Europe in the
late 1940s) introduced the circle to the Protestant David Wilkerson's \emph{The Cross and the
Switchblade}. After contact with Protestant charismatics in the area several faculty members
received the Pentecostal baptism with speaking in tongues, and by mid-February 1967, at ``the
Duquesne weekend,'' the experience had come to a wider group of students and faculty. The news
reached the University of Notre Dame, where Catholics received the baptism in home prayer meetings
helped by the Protestant Full Gospel Business Men's Fellowship. Shortly after Easter about a hundred
students, priests and faculty, chiefly from Notre Dame and Michigan State, held the first National
Catholic Pentecostal Conference there. It became annual, and by the sixth, in June 1972, the hundred
had become 11,500, including seven bishops and four hundred priests. Catholic Pentecostalism was by
then a vigorous international movement.

As the 1970s closed the Church of Rome had ventured some way from the walls of its medieval
fortress. Like its Protestant ``separated brethren'' it found the journey through the Age of
Ideologies a perilous and often uncertain pilgrimage.

\begin{keyfacts}
\keypeople{John XXIII (Angelo Roncalli); Pius XII; Paul VI; Cardinal Li\'enart; Cardinal Ottaviani;
Cardinal Alfrink; Cardinal Bea; Cardinal Suenens; Gregory Baum; David Wilkerson.}
\keydates{%
  \kd{1950}{Pius XII's encyclical \emph{Humani Generis}}
  \kd{1958}{Election of John XXIII}
  \kd{1959}{John XXIII announces a general council}
  \kd{1962--1965}{Second Vatican Council}
  \kd{1963}{Death of John XXIII; election of Paul VI}
  \kd{1965}{Vatican II's declaration on religious freedom; close of the council}
  \kd{1967}{The ``Duquesne weekend'': Catholic Charismatic Renewal begins}
  \kd{1968}{Paul VI's encyclical \emph{Humanae Vitae}}}
\keyterms{Vatican II; ``Pilgrim People''; \emph{aggiornamento}; Curia; conservatives and progressives;
Scripture and tradition; collegiality; religious liberty; ``separated brethren''; new liturgy;
\emph{Humanae Vitae}; Catholic Charismatic Renewal.}
\end{keyfacts}
